% Ad Atticum 7.6 --- headnote
% ad-atticum-07-06, 19 December 50 BC, Formiae

\textit{Cicero to Atticus, written from his Formian villa on
the fourteenth day before the Kalends of January, 19 December
50 BC (Perseus dateline:
\textit{Scr. in Formiano xiv K. Ian. a. 704 (50)}). The day
after Att 7.5. The letter is two short sections: a one-line
opening insisting that the daily ritual of correspondence be
kept up even when there is nothing new to say, and then a
single dense paragraph on the political question that is
crowding out everything else.}

\textit{Section 2 is one of the clearest statements anywhere
in the correspondence of Cicero's bind between private
judgement and public alignment with Pompey. He has canvassed
opinion and finds that almost no one favors fighting the war
that Caesar's demands are about to force. The Homeric tag
\emph{ou gar dē tode meizon epi kakon}, ``for surely this
evil that is on us now is no greater'' (from \emph{Iliad}
22.106 in Cicero's adaptation), points back to the prior
concessions --- the prorogations, the law allowing Caesar to
stand for the consulship in absence --- that have armed
Caesar against the senatorial order. The summary is the
famous formulation: ``my view will not be the same as what
I shall say.'' He will vote with Pompey because not to do so
would be wrong in him \emph{praeter ceteros}, beyond all
others; but the view he will not utter is that war must be
averted at any cost.}
